% Candidate S39. Author /root; GPT-6 Astra Pro; effort not exposed.
% 2026-11-21. Separate candidate pending independent whole-proof reading.
\begin{proposition}[S-39: Point summands cannot repair the balanced real-multiplication character]
\label{prop:semireg-rm-point-augmentation}
Let $(X,\theta,F)$ be a real-multiplication abelian fourfold of type
$(2,2)$ as in S-22. Put
\[
 A_i=\theta_i^2/2,\qquad C=A_1+A_2,\qquad
 \nu=A_1A_2=\theta^4/24,
 \qquad v_q=1-qC+q^2\nu,
 \qquad q\in\mathbf Q_{>0}.
\]
Let $G$ be an effective finite coherent translation--tensor action as in
S-21. Suppose that in the equivariant perfect category there is a direct sum
\[
 \widetilde E=\widetilde F\oplus\widetilde Q,
 \qquad e^B\operatorname{ch}(E)=Nv_q,
 \qquad N=\operatorname{rk}(E)>0,
 \tag{S39.1}
\]
where the underlying $Q$ is a nonzero finite direct sum of shifts of
skyscraper sheaves at closed points, with arbitrary positive multiplicities.
Then
\[
 \operatorname{ext}^2_G(E,E)\ge22,
 \qquad \dim\ker\sigma^2_{\widetilde E}\ge2.
 \tag{S39.2}
\]
Neither simplicity nor gluability of the lift is required. The hypothesis
is an actual equivariant direct sum; it makes no assertion about arbitrary
nonsplit point extensions or point modifications.

More precisely, write $\operatorname{ch}(Q)=\delta\nu$, where
$\delta\in\mathbf Q$ is allowed to vanish. The lower bounds for the
kernel are respectively $3$ if $\delta=0$, $2$ if
$\delta=Nq^2$, and $4$ in every other case.
\end{proposition}
\begin{proof}
We first compute the full degree-two contraction rank for a general
top-degree coefficient. For $b\in\mathbf Q$, consider the normalized
Hodge character
\[
 w_b=N-Nq(A_1+A_2)+bA_1A_2.
\]
Its rank is unchanged by dividing by $N$. Put $c=b/N$. The two
real-multiplication eigenspaces split the complex vector spaces
\[
 HT^1(X)=V_1\oplus V_2,
 \qquad V_i=T_i\oplus U_i,
 \qquad \dim T_i=\dim U_i=2,
\]
where $T_i$ acts by contraction in the holomorphic variables and $U_i$
acts by exterior multiplication in the antiholomorphic variables. This is
an exterior-algebra splitting of invariant forms, and does not assert a
product decomposition of $X$ into abelian surfaces.

Choose nonzero antiholomorphic volume forms $\overline\Omega_i$ in the
two eigenspace blocks. On $\bigwedge^2V_1$, the image is the span of
\[
 \overline\Omega_1(1-qA_2),\qquad
 \overline\Omega_1(-q+cA_2).
 \tag{S39.3}
\]
Indeed, the $\bigwedge^2U_1$ operators produce the first vector;
the $\bigwedge^2T_1$ operators produce a fixed nonzero scalar multiple
of the second; and all $T_1\otimes U_1$ operators vanish, since they
would require three antiholomorphic variables in a two-dimensional block.
This span has dimension one exactly when $c=q^2$, and dimension two
otherwise. The second block gives the analogous span with indices exchanged.

The cross operators $V_1\otimes V_2$ have four groups,
\[
 U_1\otimes U_2,\quad T_1\otimes U_2,\quad
 U_1\otimes T_2,\quad T_1\otimes T_2.
\]
Their images have different blockwise holomorphic and antiholomorphic
degrees, so are independent. Each group has rank four when its scalar
coefficient, respectively $1,-q,-q,c$, is nonzero. The first three
coefficients are nonzero. These images are also in different blockwise
Hochschild degrees from (S39.3). Thus the complete contraction rank is
\[
 R(b):=\operatorname{rank}c^2_{w_b}
 =\begin{cases}
 16,&b=0,\\
 18,&b=Nq^2,\\
 20,&b\ne0,\ Nq^2.
 \end{cases}
 \tag{S39.4}
\]
The value eighteen agrees with the separately checked balanced calculation
S-22. Multiplication of the character by $e^B$ does not change these ranks:
the creation--contraction conjugation used in S-22 gives invertible changes
of the degree-two source and target.

Since $Q$ is supported at finitely many points, it has rank zero and
$e^B\operatorname{ch}(Q)=\operatorname{ch}(Q)=\delta\nu$.
Here $\delta$ is the signed total length divided by
$\int_X\nu$; it need not be an integer. Consequently
\[
 e^B\operatorname{ch}(F)
   =N-NqC+(Nq^2-\delta)\nu.
 \tag{S39.5}
\]
The identity-sector comparison and full trace--Atiyah identity from S-21
apply to the equivariant summand $\widetilde F$. The composition through
its degree-two self-Ext space has the rank (S39.4), so
\[
 \operatorname{ext}^2_G(F,F)\ge R(Nq^2-\delta)\ge16.
 \tag{S39.6}
\]
This uses no semiregularity assumption on $F$.

The point summand contributes six more independent invariant classes.
For every closed point $x$, the local Koszul calculation recalled in S-23
identifies
\[
 \operatorname{Ext}^2_X(\mathcal O_x,\mathcal O_x)
      =\bigwedge^2T_xX.
\]
Under this identification, evaluation of
$H^0(X,\bigwedge^2T_X)$ is the usual evaluation of an invariant bivector
at $x$, and is an isomorphism. The Hochschild evaluation operator on the direct sum of shifted
skyscraper sheaves is diagonal on its individual summands. It is therefore
injective as a map
\[
 H^0(X,\bigwedge^2T_X)\longrightarrow
            \operatorname{Ext}^2_X(Q,Q),
 \tag{S39.7}
\]
because $Q\ne0$. The Hochschild evaluation map is natural for the given
equivariance. Translations act trivially on invariant bivectors, and
degree-zero tensor functors act trivially on Hochschild cohomology as in
S-21. Hence the image of (S39.7) is invariant for any chosen linearization
of $Q$. Its dimension is $\binom42=6$.

The two diagonal self-Ext groups are direct summands for the actual
equivariant decomposition (S39.1). Thus, writing
$m=\operatorname{ext}^2_G(E,E)$, we obtain
\[
 m\ge R(Nq^2-\delta)+6.
 \tag{S39.8}
\]
For $h=\dim\ker\sigma^2_{\widetilde E}$, the invariant CY4 Serre
pairing makes the obstruction and full trace maps transpose, of common
rank $m-h$. Their composition has rank eighteen by (S39.1) and S-22.
The rank inequality for a composition gives
\[
 18\ge2(m-h)-m=m-2h.
\]
Equation (S39.8) yields $h\ge2$ in all cases, and gives the three stated
refinements by (S39.4). This proves the proposition.
\end{proof}

\paragraph{Check and source scope.}
The full argument is new here apart from the explicitly cited S-21
trace comparison, S-22 creation--contraction convention and S-23
skyscraper Koszul calculation. Those are used in their checked scope.
The author has read the entire saved proof; independent whole-proof
review is pending. No conclusion is drawn for nonsplit point extensions.
